On place en une même position, un émetteur $E$ et un récepteur $R$ des ondes
ultrasonores, à la distance $d=\qty{42,5}{\cm}$ d'un obstacle. Les ondes
ultrasonores qui se propagent à partir de $E$, se réfléchissent sur l'obstacle
puis sont reçues par $R$.

Un système d'acquisition informatique permet de visualiser l'onde émise~(a) et
l'onde reçue~(b). La figure~\ref{fig:oscillogramme_1} donne l'oscillogramme
obtenu.

\begin{enumerate}
    \item Déterminer la valeur du retard temporel $\tau$ entre les ondes~(a)
          et~(b).
    \item Vérifier que la valeur de la célérité de propagation dans l'air est
          $v_{\text{air}}=\qty{340}{\m\per\s}$.
    \item On répète l'expérience en utilisant le même dispositif, et l'eau comme
          milieu de propagation. On obtient avec le même système d'acquisition
          informatique l'oscillogramme représenté sur la
          figure~\ref{fig:oscillogramme_2}. Dans quel milieu (air/eau), la
          propagation des ondes ultrasonores est plus rapide~? Justifier votre
          réponse.
\end{enumerate}

\begin{minipage}{0.55\linewidth}
    \begin{figure}[H]
        \centering
        \includegraphics[width=.8\textwidth]{2025_10_31_838bdd6d133f8b78fc82g-3}
        \caption{}
        \label{fig:oscillogramme_1}
    \end{figure}
\end{minipage}
\begin{minipage}{0.43\linewidth}
    \begin{figure}[H]
        \centering
        \includegraphics[width=0.6\textwidth]{2025_10_31_838bdd6d133f8b78fc82g-3(1)}
        \caption{}
        \label{fig:oscillogramme_2}
    \end{figure}
\end{minipage}
